\section{Extraction of a localizable scalar wave and accessible products}
\label{sec:extraction}

Throughout this Section, $l$ remains the scalar \emph{vacuum wave}
constructed from the trace.  In particular, division by $M^2$ on that
vacuum wave has not defined an operator on other vectors.  We first extract
from the physical virial field a scalar wave $S$ with actual local
approximants.  Only subsequently do we construct products of such accessible
waves.  The notation $L$ in a formal insertion still refers to the
single-insertion forms of Proposition~\ref{C:single-insertion}; a physical field
$L$ will be reconstructed in Section~\ref{sec:reconstruction}.

The construction has two distinct localization steps.  We first use
the canonical two-point spaces to select the part of the virial vacuum
wave with actual time weight two.  The actual time-M\"obius unitaries
then express that part through operations in physical diamonds.  After
proving independence of the chosen time direction, we quantify these
operations by bounded approximants and construct their analytic
products.  These products have real boundaries on spacelike patches;
the remaining real-axis reconstruction is a separate task.

\subsection{Projection into canonical diamond localization}

Write $X_{\mathrm{phys}}(O)$ for the closed real vacuum space of the
physical algebra in $O$, and $X_c(O)$ for the additive canonical real
space on the one-field representation $\Pi$.  The latter is determined by
the positive scalar, conserved-vector and traceless-tensor two-point
kernels of Section~\ref{sec:foundations}.

\Needspace{6\baselineskip}
\begin{lemma}\label{E:projection-localization}
For every bounded diamond $O$,
\begin{equation}\label{E:diamond-projection}
 P_\Pi X_{\mathrm{phys}}(O)\subset X_c(O).
\end{equation}
This assertion uses conformal covariance only of the auxiliary two-point
representations.
\end{lemma}

\begin{proof}
The actual parabolic containment
\eqref{C:actual-parabolic-localization} of
Lemma~\ref{C:localization-transport}, followed by the reduction in
Theorem~\ref{C:canonical-Pi}, gives the desired projected membership
for every parabolic diamond $N$ containing $\overline O$: actual
localization is contained in the comparison localization, and the
comparison generator acts canonically on $\Pi$.
The conjugation on $\Pi$ is the primary reality operation by wedge
Bisognano--Wichmann; the same statement holds in the reverse orientation.
We explain why these containments determine the bounded diamond.

The point of the geometry is to recover a bounded diamond from the
larger parabolic regions already controlled by the null-charge
comparison.  We carry out this recovery in the canonical real spaces,
where the complementary diamond can be described explicitly.

Use the Einstein cylinder, with coordinates $(\tau,x)\in\R\times S^3$
and round distance $d$.  By a Poincar\'e transformation and a dilation,
the physical patch and the chosen diamond can be written
\[
 |\tau|+d(x,N_0)<\pi,
 \qquad O=\{\,|\tau|+d(x,N_0)<\alpha\,\},\quad 0<\alpha<\pi.
\]
On the canonical representations,
\begin{equation}\label{E:canonical-complement}
 X_c(O)'=X_c(O^\perp),\qquad
 O^\perp=\{d(x,N_0)>\alpha+|\tau|\},
\end{equation}
where prime denotes symplectic complement.  Indeed, a connected conformal
transformation carries the wedge and its opposite to this diamond pair.
Transporting wedge modular data proves the equality.  There is no
assumption here about a physical algebra's duality.

For completeness, these transports are well-defined across the required
charts.  The two-point kernels are rational functions with the finite
Jacobian and weight factors of their integral bosonic primaries.  On
strict complex point waves in sufficiently small boxes the rational
covariance identity preserves every inner product.  Their span is dense,
so it defines local unitaries; the same identity gives local composition
and hence integration to the connected cover.  On nonsingular real
patches, the differentials preserve the future cone, so smaller strict
imaginary approach cones have distributional boundary values after
transport.  On overlaps, both boundary prescriptions are continuations
of the same rational complex kernel.  In strict spacelike separation
the two prescriptions agree.  Additivity therefore transports the real
smear spaces, even when the opposite diamond meets several charts.

Here is an explicit covering of the complement needed in
\eqref{E:canonical-complement}.  Let a target have $\tau\geq0$ and
$\theta=d(x,N_0)>\alpha+\tau$.  Choose
$\alpha<\beta<\theta-\tau$ and take the past tip of $N$ to be
$(-\beta,N_0)$.  Its future tip is $(\pi-r,q)$, where
$d(q,N_0)=r<(\pi-\alpha)/2$; it lies on the future null boundary of the
physical patch.  This parabolic diamond contains $\overline O$ with
strict margins.  The target is outside the causal future of its past tip
because $\theta>\beta+\tau$.  To put it outside the causal past of the
future tip, it suffices to choose $q$ so that
\[
 d(x,q)+r>\pi-\tau.
\]
At fixed $r$, the maximum left side is
$\min(\theta+2r,2\pi-\theta)$.  As $r$ increases to
$(\pi-\alpha)/2$, this tends to
$\min(\theta+\pi-\alpha,2\pi-\theta)$, which exceeds
$\pi-\tau$, except when $(\tau,x)=(0,-N_0)$.  Time reversal treats
$\tau<0$.  Consequently the spacelike shadows of these $N$ cover
$O^\perp$ except possibly its central antipode.

If a vector symplectically annihilates all these shadows, its pairing
with canonical real point waves vanishes on their union.  The resulting
distribution cannot be supported only at the omitted point: locally its
singular covectors are causal, from the positive-energy tube and its
conjugate, whereas a nonzero point-supported distribution has all
nonzero covector directions in its wavefront set.  It therefore
annihilates $X_c(O^\perp)$.  Taking symplectic complements in
\eqref{E:canonical-complement} proves
\eqref{E:diamond-projection}.
\end{proof}

\subsection{Timelike restrictions of the original fields}

The time-weight argument uses point waves restricted to a timelike
line.  Later, bounded approximation will also require uniform moments
of the original source field as a transverse smearing shrinks to that
line.  Both facts follow from the same ordered-spectrum estimate.

\Needspace{6\baselineskip}
\begin{lemma}\label{E:timelike-restriction}
Every fixed finite original vacuum distribution has a restriction to
parallel timelike lines, continuous in joint Schwartz temporal tests.
The restrictions of its fixed transverse derivatives exist as well.
Transverse mollification converges in these distributions, and in the
one-field vacuum and adjoint-vacuum vectors, with error $O(\delta)$
at width $\delta$.  Each bound uses finitely many temporal test
seminorms, uniformly on compact families of line positions.  In
particular every fixed source-field polynomial moment is uniformly
bounded under this mollification.  For a Hermitian field, a real line
test has its vacuum vector in $X_{\mathrm{phys}}(O)$ whenever the
compact line support is contained in the double cone $O$.
\end{lemma}

\begin{proof}
Choose the line direction as the time axis and write the individual
Fourier variables as $k=(\omega,\boldsymbol k)$.  On the ordered
spectral support $\Gamma_n$, positive energy of the partial sums gives
$\sum_i|\boldsymbol k_i|\leq C_n\sum_i|\omega_i|$.
Choose a smooth cutoff $\chi=1$ near $\Gamma_n$, supported where
$|\boldsymbol k|\leq C'_n(1+|\omega|)$, with polynomially bounded
derivatives.  For a joint Schwartz temporal test $h$,
$\chi(k)\widehat h(\omega)$ is therefore Schwartz in all variables,
with each Schwartz seminorm bounded by finitely many temporal ones.
Compact families of spatial translations and fixed transverse
derivatives introduce only polynomial factors and obey the same bound.

For the tempered Fourier distribution $\widehat W_n$, fix a finite
Schwartz seminorm $p_m$ controlling its action.  Transverse mollifier
multipliers are uniformly bounded in the required weighted derivative
seminorms, and
\[
 |\partial_{\boldsymbol k}^{\gamma}
       (\widehat\rho(\delta\boldsymbol k)-1)|
       \leq C_\gamma\delta(1+|\boldsymbol k|),\qquad 0<\delta\leq1.
\]
After multiplication by $\chi\widehat h$, this gives an $O(\delta)$
bound in $p_m$, with finitely many additional temporal seminorms.
It constructs the restriction and proves the stated distributional
convergence.  For the squared one-field vacuum error, the doubled
correlation has two mollifier differences; its seminorm is
$O(\delta^2)$, giving vector error $O(\delta)$.  The adjoint and all
fixed alternating source-field words are covered by the same original
vacuum distributions.  Thus every fixed moment required below is
uniformly bounded.  This argument tests tempered distributions; it
assumes no ordinary Fourier density for an $n$-point function.

For a real line test of a Hermitian field, sufficiently thin compact
transverse mollifications are supported in $O$.  Their vacuum vectors
belong to its closed real space by affiliation.  Their Hilbert limit
has the same membership.  Translation and Lorentz covariance give the
assertion for all the lines under consideration.
\end{proof}

\subsection{Actual and auxiliary time weights}

Fix a future timelike unit $e$ and put $E=P_e>0$.  Use the common affine
realization $L^2(\R_+,\dd E/E;\mathcal K)$, where translations multiply
by $e^{itE}$ and dilations rescale $E$.  In a lowest-weight module of
weight $h>0$, the point wave is $E^h e^{izE}v$.  Here $h$ is a
scalar time lowest weight, unrelated to the internal parabolic operator
$h$ of Proposition~\ref{F:parabolic-generator} or the radial comparison
operators $h_\nu$.  Theorem~\ref{C:time-mobius} and
Lemma~\ref{C:mobius-real-spaces} constructed the \emph{actual}
time-M\"obius action and its interval real spaces.  The auxiliary
canonical representation has the same affine convention, with its own
primary weights $s$.  Write $H$ for the actual lowest-weight operator
on the affine multiplicity space.  We will compare interval real
spaces as well as affine actions: the affine actions alone do not
distinguish lowest weights.

\Needspace{6\baselineskip}
\begin{lemma}[Interval Paley regularity and weight selection]
\label{E:paley-selection}
Let $X_s(I)$ be the closed real interval space of a primary of weight
$s>0$.  Dividing a wave in $X_s(I)$ by $E^s$ gives an entire Fourier
transform; evaluation and each derivative on a compact complex-frequency
set are continuous in its real-space Hilbert norm.  The bounds are
uniform on compact positive weight bands.  In particular:
\begin{enumerate}
\item Suppose a bounded affine intertwiner $j_s$ from the actual
representation to a weight-$s$ primary representation satisfies
$j_sX_{\mathrm{act}}(I)\subset X_s(I)$ for bounded intervals $I$.
In the common affine realization write $j_s=\id\otimes T_s$.
Then
\[
 T_s\mathbf1_{\R\setminus(s+\mathbb Z_{\geq0})}(H)=0.
\]
\item A physically real localized vacuum point wave of scale dimension
$d>0$ has actual weights only among $h=d-j>0$,
$j\in\mathbb Z_{\geq0}$.
\end{enumerate}
\end{lemma}

\begin{proof}
For real $f\in C_c^\infty(I)$ the primary norm, up to a positive
constant, is
\[
 \int_0^\infty E^{2s-1}|\widehat f(E)|^2\,\dd E.
\]
Reality determines negative frequencies.  To evaluate
$\partial_z^k\widehat f(z)$, choose a compactly supported smooth test
which agrees on $I$ with $(it)^k e^{izt}$.  Modify it outside $I$ to
annihilate its first $m$ Fourier jets at zero, with $m>s-1$.
Such a modification exists: finitely many disjoint small bumps outside
$I$, with distinct centers, have a moment matrix arbitrarily close to
a nonsingular Vandermonde matrix.  Cauchy--Schwarz with dual weight
$|E|^{1-2s}$ now bounds the pairing.  Its high-frequency integral
converges by smoothness, and its low-frequency integral converges by
the vanished jets.  The construction and all bounds are uniform on
compact sets of $z$ and compact bands of $s$.  They extend to the real
Hilbert closure.  The extended evaluations satisfy the
Cauchy--Riemann equations and give the asserted entire function.

The regularity at zero frequency now detects the allowed weights.
In the common affine realization, changing the primary weight changes
only a power of $E$.  Entire continuation permits precisely the
nonnegative integral powers of that variable.

An intertwiner commuting with $E$ and its rescalings is an
energy-independent bounded multiplicity map $T_s$.  For a real
multiplicity section $v$ in a compact $H$-band, projection of the actual
localized wave has quotient
\begin{equation}\label{E:weight-multiplier}
 \widehat f(E)T_s E^{H-s}v.
\end{equation}
The assumed real-space inclusion and the preceding continuity make
this holomorphic at $E=0$.  Choose $\widehat f(0)\ne0$.  On a compact
band $T_s\exp((H-s)z)v$ is an entire function of $z=\log E$.
Single-valued continuation in $E$ implies
\[
 T_s\bigl(e^{2\pi i(H-s)}-\id\bigr)v=0.
\]
Complexifying real sections and applying the spectral theorem shows
that $T_s$ is supported on $h-s\in\mathbb Z$.  Only finitely many
such values occur on the band.  Their finite Laurent expansion in
\eqref{E:weight-multiplier} has no pole, so its negative powers vanish.
Exhausting the positive weight axis proves the first assertion.

For the second, write a dimension-$d$ point wave as $E^d u$ in the
actual affine representation.  Project onto compact $H$-bands, which
commute with interval Tomita data.  In each fiber division by its
primary power gives $\widehat f(E)E^{d-h}u_h$.  The same argument forces
$d-h\in\mathbb Z_{\geq0}$.  There is no translation-invariant component
of a positive-dimension point wave: its zero-momentum vacuum amplitude
would have to transform by a nonunit-modulus dilation character.
For original fields, the required timelike restrictions and real
localization are supplied by Lemma~\ref{E:timelike-restriction}.
The same assertion for a subsequently constructed wave uses its stated
Hilbert-valued restriction and physical real localization.
\end{proof}

The norm-continuous evaluations above concern the closed real space.
For a complex localized operator vector, use its Hermitian real and
imaginary parts; the bound then involves the sum of the vacuum and
adjoint-vacuum norms.  No Hilbert-norm evaluation bound on an arbitrary
completion of $X_s(I)+iX_s(I)$ is being used.

Here the intertwiners are specific canonical projections.  Project by
$P_\Pi$, then onto an auxiliary time-primary summand of weight $s$.
Lemma~\ref{E:projection-localization} and the canonical interval
Tomita data give the required real-space inclusion for this map $j_s$.
The same construction applies when the canonical comparison space is
enlarged below.  Taking adjoints of the selection identity gives
\[
 \mathbf1_{\R\setminus(s+\mathbb Z_{\geq0})}(H)T_s^*=0.
\]
Write $l(0)=E^2u_l$.  Its auxiliary temporal primary has weight two,
so its actual affine label is $u_l=T_2^*v_l$ for its canonical label
$v_l$; it has actual weights in $2+\mathbb Z_{\geq0}$.
Likewise $J_e(0)=E^3u_{J_e}$ has
$u_{J_e}=T_3^*v_{J_e}$ and actual weights in
$3+\mathbb Z_{\geq0}$.  These are adjoint support statements, not
commutation of actual and auxiliary weight projections.
The temporal virial wave has $E^3$-label $iu_l+u_{J_e}$.
Its physical real localization and dimension three bound its actual
weights above by three.  It therefore has only weights two and three.
The weight-two label is $iu_{S_e}$, where
\begin{equation}\label{E:definition-Se}
 u_{S_e}=P^{\mathrm{act}}_{h=2}u_l,
 \qquad S_e(0)=E^2u_{S_e}.
\end{equation}
This is a vacuum-wave definition.

The derivative of a primary of weight two has an inhomogeneous
transformation term.  Differentiating its primary transformation, and
adding the weight-three part of $V_e$, gives for a small real
non-affine time-M\"obius map $g$
\begin{equation}\label{E:finite-difference}
 U_e(g)V_e(ze)\Omega
  =(g'(z))^3V_e(g(z)e)\Omega
     +2g'(z)g''(z)S_e(g(z)e).
\end{equation}
Choose an interval on which $g'g''\ne0$.  After changing variables and
multiplying by smooth real coefficients, this equation expresses any
smearing of $S_e$ there as a difference of a physical virial smearing
and an actual modular transform of one.  Consequently, for a real test
supported in a line interval whose diamond lies in $O$,
$S_e(f)\in X_{\mathrm{phys}}(O)$.  Complex linear extension gives its
adjoint-wave reality relation; this is the meaning of a Hermitian wave
here.  It is even and physically real localized on all translated lines.
Translation
covariance defines it as a spacetime Hilbert-valued distribution.  For
example $\|e^{-\epsilon P_e}S_e(0)\|=O(\epsilon^{-2})$,
and $|\boldsymbol P|\leq P_e$ supplies all required spacetime test
bounds.  The family is Lorentz covariant in $e$, a rotational scalar
about $e$, and homogeneous of dimension two.

\Needspace{9\baselineskip}
\subsection{Independence of the time direction}

\begin{proposition}\label{E:accessible-scalar}
The waves $S_e$ agree for all future timelike unit vectors $e$.
Their common value $S$ is a purely timelike Lorentz scalar of dimension
two, with physical real localization.  It has the canonical null
special-conformal action.  It can be included in the auxiliary
representation and in Lemma~\ref{E:projection-localization}.
\end{proposition}

\begin{proof}
There are three issues to settle: the timelike spectral densities,
the possible massless component, and the action of the actual null
generators on the resulting scalar.  The projection identity
\eqref{E:definition-Se} will link the first two issues: comparison on
one energy slice fixes both the normalization and the dependence on
the time direction.

We first prove a two-point claim: the timelike entries of the joint
spectral measure of the family $S_e$ and $l$ are constant multiples of
the ordinary measure $\dd^4p$ on the future timelike cone.  No density
is assumed in this claim.  Write $\mu_{ee'}$ for a joint spectral
entry, including $l$ as an additional label.  The corresponding
commutator distribution is
\[
 C_{ee'}(x)=\int e^{ip\cdot x}\,\dd\mu_{ee'}(p)
              -\int e^{-ip\cdot x}\,\dd\mu_{e'e}(p).
\]
Every invariant-mass projection commutes with wedge Tomita data.
Wedge separation and reality therefore give causal support for this
commutator after projection to any compact mass interval.  For $l$,
the real wedge data are the canonical scalar data already identified
in Section~\ref{sec:foundations}; physical point-field membership is
not needed.

Homogeneity disintegrates the timelike measure as
$M^3\dd M\,\dd\nu(p/M)$, with a polynomially bounded measure on the
unit hyperboloid.  Testing arbitrary mass intervals shows that the
shell commutator has causal support for almost every mass; rescaling
then gives the same statement for the reference shell.  This is a
disintegration of distributions, not an assumed field on a fixed
mass shell.  The shell commutator solves the Klein--Gordon equation.
Its displacement and normal-derivative Cauchy data at zero time are
supported at the spatial origin, and hence are finite sums of delta
derivatives.  Fourier solution of this Cauchy problem expresses its
shell measure as the ordinary commutator measure times a polynomial
in $p/M$, modulo $p^2=M^2$.  In particular this proves absolute
continuity rather than presupposing it.

For a self-pair the shell polynomial is even under $p\mapsto-p$.
Take its representative of least degree modulo the shell relation.
If its highest degree is $k$, that homogeneous term is nonzero on an
open set of null directions; otherwise divisibility by $p^2$ would
reduce the degree.  At fixed nonzero energy and $M\downarrow0$ it
contributes order $M^{-k}$ to the density with respect to $\dd^4p$.
The transverse mass integration is $M\dd M$.  Positivity and local
finiteness therefore require $k<2$.  Evenness leaves only a constant.
Positivity of the matrix measure bounds each cross density by the
geometric mean of the self densities.  A polynomial on the unit
hyperboloid with such a global bound is constant by the same highest
homogeneous term argument.  All timelike density entries are therefore
constant, proving the two-point claim.  Covariance makes the self
constant of $S_e$ independent of $e$.

There is no zero-momentum part.  We must also exclude a massless part;
positivity alone would allow one.  Let $\rho=\sqrt{-W_\mu W^\mu}$ be
the massless Pauli--Lubanski length and project the entire Lorentz family
onto a compact band $0<a\leq\rho\leq b<\infty$.  This projection
commutes with wedge modular data.  Its commutators are causal massless
solutions.  The preceding point-supported Cauchy-data argument makes
them derivatives of the massless commutator, hence supported on the
light cone in four dimensions.

On the complex cyclic span of this projected family, let $X_+$ and
$X_-$ be the real spaces generated in the future and past cones.
Translation-spectrum uniqueness makes both spaces complex cyclic:
a vector orthogonal to smearings on an open cone is orthogonal to all
translates by tube uniqueness.  Huygens support gives
$X_-\subset X_+'$, so $X_+$ is standard.  Lorentz transformations
preserve it and hence commute with its modular group.  Every future
translation maps it into itself.  The positive-generator modular
translation relation therefore scales all $P_\mu$ by one common
nontrivial exponential under that modular group.  Because the Lorentz
generators commute with the group, $\rho$ scales by the same factor.
No nonzero spectrum contained in $[a,b]$ is invariant under all these
rescalings.  The projected space is zero.  Exhausting positive bands
excludes nonzero Pauli--Lubanski length.

The remaining massless representations have trivial little-group
translations and decompose into helicities.  A rotational scalar about
$e$ has only helicity zero, since a spatial rotation fixing both $e$ and
a massless momentum acts trivially on its wave.  Use the scalar
massless invariant measure.  The possible internal part of dilations
is a momentum-independent unitary group $W$ on scalar multiplicities.
Lorentz covariance and dimension two give the amplitude
\begin{equation}\label{E:massless-amplitude}
 S_e(p)=(p\cdot e)W(\log(p\cdot e))\beta.
\end{equation}
For two directions the commutator argument makes the cross density a
homogeneous quadratic shell polynomial.  A quadratic scalar invariant
under rotations fixing $e,e'$ is, on the null shell, a linear
combination of $(p\cdot e)^2$, $(p\cdot e)(p\cdot e')$ and
$(p\cdot e')^2$.  It follows that
\[
 \varphi(u)=\ip{\beta}{W(\log u)\beta},\qquad
 u=\frac{p\cdot e'}{p\cdot e},
\]
is a linear combination of $u^{-1},1,u$ on the attainable interval.
As the relative rapidity grows these intervals exhaust $(0,\infty)$;
the coefficients agree on overlaps.  Since
$|\varphi(u)|\leq\|\beta\|^2$, both nonconstant coefficients vanish.
Thus the massless cross density is
$c(p\cdot e)(p\cdot e')$, with $c=\|\beta\|^2\geq0$.

The only massless possibility left is the scalar density just
displayed.  Its growth with the relative rapidity of $e$ and $e'$ is
incompatible with the bounded timelike cross densities.  The next
energy-slice calculation makes this comparison with its exact
normalization.

Now apply the defining projection in \eqref{E:definition-Se} to
$F=S_{e'}$.  A dimension-two wave with physical real localization has actual weights at most
two; $u_l$ has actual weights $2,3,\ldots$.  In the multiplicity fiber
at $P_e=1$,
\begin{equation}\label{E:projection-pairing}
 \ip{u_F}{u_{S_e}}=\ip{u_F}{u_l}.
\end{equation}
Let $B(e',e)$ and $D(e')$ denote the constant timelike density
entries for $(S_{e'},S_e)$ and $(S_{e'},l)$, respectively.  Write
$b=B(e,e)$, which is independent of $e$ by covariance, and let
$\rho_{ll}$ be the self density of $l$.  With timelike measure
$\dd^4p$ and massless measure $\theta(p^0)\delta(p^2)\dd^4p$,
\eqref{E:projection-pairing} reads explicitly
\begin{equation}\label{E:energy-slice-normalization}
 \frac{4\pi}{3}B(e',e)+2\pi c(e\cdot e')
   =\frac{4\pi}{3}D(e').
\end{equation}
Indeed the timelike slice $p\cdot e=1$ is the unit spatial ball
in the rest frame of $e$, of volume $4\pi/3$.  Lorentz invariance
of $\dd^4p$ makes this normalization independent of the slicing
direction.  On the massless slice the angular integral is
$\frac c2\int_{S^2}(e'^0-\boldsymbol e'\cdot\boldsymbol n)\dd\Omega
=2\pi c(e\cdot e')$.  Positivity gives
$|B(e',e)|\leq b$ and $|D(e')|\leq\sqrt{b\rho_{ll}}$ uniformly.
Arbitrarily large relative rapidity therefore forces $c=0$.
Setting $e=e'$ in \eqref{E:energy-slice-normalization} now gives
$D(e')=b$, and hence $B(e',e)=b$ for all pairs.  Thus
$\|S_e(f)-S_{e'}(f)\|^2=0$ for every test $f$.
It also shows that the constant two-point density of $S$ with $l-S$
is zero.  This is orthogonality of the scalar wave channels, including
all their spacetime translates; it does not assume that the actual
time-weight projection commutes with spatial translations.

Finally, physical real localization supplies the scalar order-zero
parabolic comparison for $S$.  We can also apply the proof of
Lemma~\ref{C:K-domains} to its smooth compact vacuum smearings
in the same strictly one-sided placements, and then to all their
strictly energy-damped translates.
That proof requires real localization, giving the Tomita bound after
parabolic transport, and boost smoothness of every order.  Both are
available here: localization was obtained from
\eqref{E:finite-difference}, and the scalar covariance just proved
identifies every iterated boost derivative with another compact
smooth $S$-smearing.  Its Tomita and Hardy estimates therefore give
the asserted $K_n$ domains without treating $S$ as an operator.
Symmetry against the same dense core transfers the even second
difference of \eqref{C:quadratic-charge-translation},
$K_n(s)+K_n(-s)-2K_n=4s^2a$, to these vectors as a form identity.
The $D'$ term cancels, so no $D'$ operator domain for $S$ is assumed.
The second-difference and zero-capacity argument of
Theorem~\ref{C:canonical-Pi} therefore applies: the
nonnegative comparison defect is translation-diagonal in $M^2$, and
is bounded by $\int_0^\infty M|f'(M)|^2\dd M$.  Logarithmic cutoffs
equal to one on a prescribed compact mass interval make this bound
tend to zero.  Spectral comparison then gives reduction and the exact
canonical null action.  Its wedge reality and two-point kernel are
scalar, so the proof of Lemma~\ref{E:projection-localization} includes
this enlarged canonical space as well.
\end{proof}

\Needspace{6\baselineskip}
\begin{corollary}\label{E:infinity-vanishing}
Put $\ell=l-S$.  Its insertion at imaginary infinity annihilates every
compact physical vacuum polynomial:
\begin{equation}\label{E:residual-infinity}
 \lim_{R\to\infty}R^4\ip{\ell(iRe)}{A\Omega}=0.
\end{equation}
\end{corollary}

\begin{proof}
The multiplicity label $u_\ell$ has actual weights $2+j$, $j\geq1$.
First take a Hermitian polynomial localized in a diamond on the line.
Its vacuum vector $\psi$ is in the corresponding real space.  The
scalar quotient
\[
 F(E)=E^{-2}\ip{u_\ell}{\psi(E)}
\]
is regular at zero, with continuous evaluation, by
Lemmas~\ref{E:projection-localization} and~\ref{E:paley-selection}.
For compact actual weight-band projections of $\psi$, its overlap with
$u_\ell$ has only positive integral powers $E^{h-2}$, so $F(0)=0$.
These projections preserve the interval real space and converge in
Hilbert norm; continuity of evaluation proves the same equality
without the band cutoff.  With the common $\dd E/E$ normalization,
\[
 \ip{\ell(iRe)}{\psi}=\int_0^\infty E^3e^{-RE}F(E)\,\dd E.
\]
The contribution near zero is $o(R^{-4})$, and the complementary
contribution is exponentially small by Cauchy--Schwarz.  Every compact
polynomial is a complex linear combination of Hermitian ones, and its
support fits in a line diamond.  For a compact joint polynomial
kernel, the same real-space statement is the closed Tomita-graph
limit of factorized tests, using the old joint-distribution
continuity from Section~\ref{sec:foundations}.  This proves the claim.
\end{proof}

At this point $S$ is a direction-independent scalar vacuum wave
with physical real localization, and $l-S$ has the stated decay against
physical polynomial vectors.  To multiply $S$ with other entries we
need quantitative bounded operations, including control of both their
vacuum and adjoint-vacuum errors.

\subsection{Bounded approximants for accessible waves}

Call a vacuum wave \emph{accessible} when it has the approximation
property below, including the corresponding adjoint wave.  Physical
fields and $S$ have this property.  This definition does not assert
that an accessible wave already acts on a common field domain.
In the following norm estimates, $a(f)$ denotes a vector: for an
original field it is its smeared vacuum vector, and for $S$ it is
its Hilbert-valued smearing.  The notation
$a^\dagger(\bar f)$ denotes the corresponding adjoint-vacuum vector.

\Needspace{6\baselineskip}
\begin{proposition}\label{E:bounded-approximation}
Let $a$ be a physical field or $S$, and let $f$ have support in a small
real source patch.  Given a double-cone neighborhood of its support,
there are bounded operations $A_R(f)$ in that neighborhood, for dyadic
$R\to\infty$, such that for one fixed exponent $q$
\begin{align}
 \|A_R(f)\|&\leq C R^q,\label{E:approx-norm}\\
 \|A_R(f)\Omega-a(f)\|+
 \|A_R(f)^*\Omega-a^\dagger(\bar f)\|
   &\leq C_N R^{-N}\quad(N=1,2,\ldots).
 \label{E:approx-error}
\end{align}
For each prescribed $N$, constants are uniform on balls of finitely
many test seminorms.  When patches and localization margins shrink to
size $\varepsilon$ in a fixed compact set, these constants cost at
most fixed inverse powers of $\varepsilon$.  The approximants need
not depend linearly on $f$.
\end{proposition}

\begin{proof}
For an ordinary field smearing with closed affiliated realization
$A=v|A|$, take its polar cutoff $A_R=v|A|\mathbf1_{[0,R]}(|A|)$.
It has norm at most $R$.  The tails on the vacuum and adjoint vacuum
satisfy, for every positive integer $m$,
\begin{align*}
 \|(A-A_R)\Omega\|&\leq R^{-m}\||A|^{m+1}\Omega\|,\\
 \|(A^*-A_R^*)\Omega\|&\leq R^{-m}\||A^*|^{m+1}\Omega\|.
\end{align*}
Here the absolute-value moments follow from the original realization
hypothesis together with the generated-domain action check in
Section~\ref{sec:foundations}. Both $A$ and $A^*$ preserve the dense
invariant vacuum-generated domain $\mathcal D_{\mathrm{old}}$ and agree
there with the physical field and its polynomial adjoint. Induction
therefore gives
\[
 \mathcal D_{\mathrm{old}}\subset
 \Dom (A^*A)^m\cap\Dom(AA^*)^m
 \qquad(m=1,2,\ldots),
\]
and these powers on that domain are the corresponding alternating
polynomial words.  Their vacuum norms are finite and are bounded by
finitely many seminorms of the source tests, by the joint Wightman
distributions.  Since $|A|^{2m}=(A^*A)^m$ and
$|A^*|^{2m}=(AA^*)^m$, spectral calculus also gives every lower
integer absolute-value moment used in the tail bounds.  This uses
the adjoint-extension clause of Assumption~\ref{ass:framework} and
does not require essential self-adjointness.

The invariant polynomial domain supplies these moments, uniformly
on finite test-seminorm balls; subdivision supplies the requested
local neighborhoods.  Use a Hermitian basis or adjoint-paired
cutoffs to preserve adjunction and grading.

For a nonhomogeneous original source, any comparison of shrinking
patches by dilation is made on its finitely many homogeneous
components and then summed.  For this fixed field and requested
accuracy, a maximum of the finitely many seminorm orders and powers
is sufficient.  This does not assign the source a single scaling
dimension or require a bound uniform in the set of dimensions, and
it does not change the fixed polynomial norm exponent in
\eqref{E:approx-norm}.

For $S$ the order of construction matters.  We truncate the
ordinary source field before applying the actual modular unitary,
and then add the resulting bounded operations.  We will estimate
the sum by its absolute quadrature weights, so no domain or moment
estimate for an unbounded sum of modularly transformed fields enters.

For $S$, use \eqref{E:finite-difference} on parallel lines.  After
undoing $g$ in the test, a line smearing of $S$ is a fixed linear
combination of a line virial smearing and a modularly transformed
line virial smearing, with smooth real coefficients.  By
Lemma~\ref{E:timelike-restriction}, transverse approximation at width
$\delta_R\leq e^{-R}$ has vacuum error $O(\delta_R)$, and every
fixed source-field moment is bounded uniformly as
$\delta_R\downarrow0$.  Its adjoint data obey the same estimates.

Polar truncate each such source field at $R$ and apply the actual
M\"obius unitary.  It fixes the vacuum and transports the bounded
operation into the destination diamond, so both estimates in
\eqref{E:approx-error} survive.  A transverse quadrature of mesh
at most $e^{-R}$ approximates a full spacetime smearing.  Although
its number of terms grows, the sum of absolute quadrature weights
is bounded.  Hence the sum of the operator norms is $O(R)$, and the
sum of the vacuum errors remains rapidly small.  No moment bound
for a sum of noncommuting transformed fields is being assumed.

All the operations just described use finitely many derivatives for
each requested accuracy.  Rescaling tests and using diamonds with
margins proportional to $\varepsilon$ therefore introduces only
inverse powers of $\varepsilon$.  The fixed non-affine map can be
chosen nonsingular on the entire compact parameter patch, with
$g'g''$ bounded away from zero.  Take
$\delta_R\leq c\varepsilon e^{-R}$ and sufficiently fine quadrature
to retain these margins.  Subdivision completes the construction
for an arbitrary compact source support.
\end{proof}

The norm exponent in \eqref{E:approx-norm} is fixed before an error
order $N$ is requested.  Increasing $N$ changes the finite collection
of moments and test seminorms used to estimate the error, but does not
change that exponent.  This distinction allows rapid vacuum errors
to dominate the finitely many polynomially growing factors in a
product.

\subsection{Holomorphic products, ties, and unsmearing}

For $p$ accessible entries, use increasing imaginary translations:
\begin{equation}\label{E:ordered-tube}
 \operatorname{Im}z_1\in V_+,
 \qquad\operatorname{Im}(z_{j+1}-z_j)\in V_+
 \quad(1\leq j<p).
\end{equation}
The first inequality is needed for vector products.  With tests fixed
in sufficiently small patches, form actual bounded vector products
\begin{equation}\label{E:bounded-product}
 F_R(z)=U(z_1)A_{1R}U(z_2-z_1)A_{2R}\cdots
           U(z_p-z_{p-1})A_{pR}\Omega.
\end{equation}
They are holomorphic in \eqref{E:ordered-tube}, with norm at most a
fixed power of $R$.
Their scalar vacuum expectations have the larger tube determined only
by the successive imaginary gaps: vacuum invariance removes the first
factor $U(z_1)$.  Equivalently, add a common future imaginary translation
to put all the arguments in the vector tube, then pair with $\Omega$.
The scalar answer is independent of this common translation.  We use
$W$ for its limit and retain $\Psi$ for the vector limit.

\Needspace{6\baselineskip}
\begin{proposition}\label{E:accessible-products}
The vectors \eqref{E:bounded-product} have a unique, locally uniform
holomorphic limit.  It is multilinear in the accessible waves and has
Hilbert-valued distributional boundaries at mutually strictly
spacelike real placements.  These boundaries are graded symmetric.
The limits extend to common holomorphic neighborhoods of imaginary
ties: successive blocks may have equal imaginary translation when
real placements within each block are pairwise spacelike and all
inequalities between blocks remain strict.  The limits unsmear to
holomorphic point vectors, denoted $\Psi$, depending only on the total
complex insertion arguments.  For physical entries alone they agree
with the usual Wightman vectors.
For any two finite accessible words, with each word in its strict
vector tube, their scalar functions satisfy
\begin{equation}\label{E:accessible-splitting}
 \begin{split}
 &\big\langle\Psi[A_1(z_1)\cdots A_p(z_p)],
                  \Psi[B_1(w_1)\cdots B_q(w_q)]\big\rangle\\
 &\quad=W\bigl(A_p^\dagger(\bar z_p),\ldots,
               A_1^\dagger(\bar z_1),B_1(w_1),\ldots,B_q(w_q)\bigr).
 \end{split}
\end{equation}
An empty word denotes $\Omega$.  Source tests on the bra word are
conjugated and reversed, together with their field labels.  The identity
holds for distinct words and their finite linear combinations at every
finite total length.  It asserts analytic Hilbert splittings and their
spacelike real boundaries, without yet asserting a distributional
boundary at every real configuration.
\end{proposition}

\begin{proof}
We first prove convergence where locality lets us estimate the
whole product by vacuum errors.  Holomorphic propagation then carries
that convergence into the ordered tube.  Finally, bounded-product
continuity at a spacelike edge gives the common neighborhoods needed
for ties; translation compatibility removes the source smearing.

At real translations making all approximation patches mutually
spacelike, telescope two bounded products, whether at $R,2R$ or
from different approximation choices.  Strong graded locality moves
each difference to the vacuum.  Its vacuum error is rapidly small
by \eqref{E:approx-error}, while all other factors have polynomial
norm growth by \eqref{E:approx-norm}.  The product difference is
therefore $O_N(R^{-N})$ for every $N$, uniformly on smaller placement
patches and appropriate test-seminorm balls.  This proves independence,
graded symmetry and multilinearity at these real placements.

We record the propagation estimate, since vacuum convergence alone
would not justify the products elsewhere.  For a holomorphic vector
difference $G_R$ bounded by $CR^q$, restrict to a complex line
$z=x+\zeta\eta$ with $\eta$ in the ordered cone.  A real interval $I$
can be chosen inside a good spacelike-placement patch.  The
subharmonic logarithm of its norm and harmonic measure give
\begin{equation}\label{E:harmonic-propagation}
 \log\|G_R(x+\zeta\eta)\|
 \leq\omega(\zeta,I)(\log C_N-N\log R)
 +(1-\omega(\zeta,I))(\log C+q\log R).
\end{equation}
On compact subsets reached this way $\omega$ is bounded below by a
positive number.  Interior three-circle chains propagate the same
conclusion to every fixed compact subset of the connected tube.
As $N$ is arbitrary, convergence is rapid there.  Above a smaller
good real patch it is rapid uniformly up to the boundary, because
$\omega$ has a uniform positive lower bound for bounded approach
directions.  Summing the dyadic differences proves existence.

For ties, first keep the common translation of each block and its
spectators strictly inside their allowed complex domains.  The
relative translations within a block have wedges of opposite
imaginary-gap cones for its order and its reversal.  At their common
real edge the bounded products agree, with the graded sign, by
strong locality.  Each bounded product has a norm-continuous
extension to that edge: the real translation groups and the
positive-energy contraction semigroups are strongly continuous,
and only finitely many bounded factors occur.  On fixed smaller
wedges and their common edge its norm is $O(R^q)$.
The continuous-edge form of the local edge-of-the-wedge theorem
therefore applies directly; its continuation neighborhood depends
only on the wedges and the edge patch, and its interior bounds are
controlled by the supremum bounds on those wedges and their edge;
see \cite[Theorems~1$'$ and~2]{Browder1963}.

To make uniformity explicit, apply the theorem after scalarization
to the Hilbert-valued family $(R^{-q-1}F_R)_R$ in the
square-summable direct sum, with dyadic $R$.  Its common edge is
norm continuous: finitely many components are continuous and the
remaining squared norms are bounded uniformly by a convergent tail
of $\sum_R R^{-2}$.  Scalarization gives a bounded holomorphic
Hilbert-valued continuation on one common neighborhood, and each
unweighted component there has a bound $O(R^{q+1})$.  The same
supremum-bound argument is uniform on the prescribed test-seminorm
balls; no estimate from boundary distribution seminorms alone is
being asserted.  Interior
propagation from the strict tube now gives rapid convergence on
compact subsets of these neighborhoods.  Adjacent permutations
within a block are treated on overlapping smaller neighborhoods,
where uniqueness glues the continuations.

By uniqueness from the spacelike real patch, multilinearity persists
throughout these domains.  The seminorm bounds and the test-function
kernel theorem give joint distributions in the source variables,
holomorphic in the translation variables.  Translate a source test
within its patch.  At real spacelike placements this is the same as
translating its insertion argument; the equality involves only the
vacuum wave and hence survives all approximation choices.  Analytic
uniqueness extends this equality.  If $x$ denotes source variables,
the kernel in coordinates
$(x,\operatorname{Re}z+x,\operatorname{Im}z)$ is therefore independent
of its first variables in the distributional sense.  It satisfies
the Cauchy--Riemann equations in the remaining variables.  Thus it is
a holomorphic Hilbert-valued function of $x+z$, rather than merely a
distribution in $x$.  This proves unsmearing.  The same identity glues
source patches and their translations.

For the split identity first keep all source tests smeared, and use
the same bounded approximants on the two sides, taking the adjoints of
the bra approximants in the scalar word.  This is an admissible choice
because both vacuum and adjoint-vacuum errors were controlled.  Bounded
adjunction gives \eqref{E:accessible-splitting} before taking limits.
The displayed scalar word lies in its ordered scalar tube: its reversed
bra gaps and ket gaps are future directed, and its central gap has
imaginary part $\operatorname{Im}w_1+\operatorname{Im}z_1\in V_+$.
Local uniform convergence of the two vector factors and of the scalar
word therefore proves the identity there.  The preceding unsmearing
then gives the point formula.  Conversely, at any cut between two
successive entries of an ordered scalar tuple, choose a common
imaginary shift strictly between their imaginary parts.  Subtracting
that shift puts the reversed-adjoint prefix and the suffix in the two
strict vector tubes of \eqref{E:accessible-splitting}.  Thus the formula
supplies the analytic Hilbert splitting of every ordered scalar tuple
at every cut.  Finite linear combinations retain all
cross terms.  At mutually spacelike real placements, the same identity
follows from the Hilbert-norm limits already proved on those patches;
no unrestricted real-axis limit is needed.

For physical fields, polar truncations converge on the polynomial
domain as well as on the vacuum.  At mutually spacelike real
placements their product limit is consequently the ordinary smeared
product.  For a joint source kernel, finite tensor-test approximation
and the joint distribution estimates extend this equality; the
kernel itself need not have finite tensor rank.  Tube uniqueness
identifies their point vectors everywhere.
Finally the estimates above give distributional boundary convergence
locally uniformly on the good real patches.  In any closed strict
subcone of approach, point bounds near such an edge are polynomial
in inverse distance: test the real variables on balls at that scale
and use the holomorphic mean-value estimate on smaller complex balls.
\end{proof}

\subsection{A boundary pullback rule}

The preceding products will be composed with analytic coordinate
changes before their real limits are taken.  The following formulation
specifies the required boundary operation.

\Needspace{6\baselineskip}
\begin{lemma}\label{E:boundary-calculus}
Let a Hilbert-valued holomorphic function have a local distributional
boundary $u$ and polynomial distance bounds in closed strict approach
subcones.  Let $\Phi$ be a real analytic local diffeomorphism whose
complexification takes a smaller wedge strictly into such approach
directions.  Then its composition with $\Phi$ has boundary
$\Phi_{\R}^{*}u$ on that smaller wedge.
\end{lemma}

\begin{proof}
Shrink the real patch and the approach cone so that the image of a
small imaginary translate of size $\delta$ is a graph at distance
comparable to $\delta$ inside one fixed strict convex subcone.
Interpolate that graph to a straight imaginary translate there.
Suppose the holomorphic function is $O(\delta^{-m})$ on this
interpolation.  Extend a compact smooth test almost analytically to
order $N$: its $\bar\partial$ is $O(\delta^N)$, uniformly with the
finitely many derivatives required on the patch.  Stokes' formula
for the holomorphic volume form compares the two integrals.  Their
error is bounded by a constant times $\delta^{N+1-m}$; take
$N>m+1$.

There is a minor test-function issue in this comparison.  The desired
coefficient on the moving graph is a fixed smooth real test, whereas
an almost-analytic extension evaluated on that graph has nonzero
imaginary Taylor terms.  Include the holomorphic Jacobian, and choose
the real test on the straight surface by inverting these Taylor terms
through order $N$.  More explicitly, the graph evaluation operator
on test jets is $\id+\delta B_1+\cdots+\delta^N B_N$, where each
$B_j$ is a differential operator with smooth coefficients.  It has a
unique inverse modulo $\delta^{N+1}$, obtained successively in each
power of $\delta$.  This is the negative-imaginary-displacement
Taylor correction.  Its corrected real tests converge in every
smooth seminorm to the ordinary real change-of-variables test; the
remaining graph error is $O(\delta^{N+1-m})$.  The boundary limit on
the straight surface is therefore precisely the distributional real
pullback.  Pairing with arbitrary Hilbert vectors proves the weak
statement, and the uniform estimates prove the asserted
Hilbert-valued distributional convergence.
\end{proof}
